\section{Exponent versus binomial index}\label{sec:position}

Two positivity statements now sit on the same one-parameter family. The
transfer's rational factor exchanges complete monotonicity for innerness at
$m=2$ (Remark~\ref{rem:trade}); the point count loses complete monotonicity on
$m\in(1,2)$ (Corollary~\ref{cor:nolattice}). The intervals look the same. This
section shows that they are not the same statement, that the agreement is forced
by arithmetic of a very cheap kind, and what does separate the two sides.

\subsection{The transfer has exactly one positivity threshold}\label{sec:onethreshold}

\begin{proposition}[One threshold, at $m=2$]\label{prop:onethreshold}
In the decomposition $K_\omega=R_\omega K^\Gamma_\omega K^\zeta_\omega$:
\begin{itemize}
\item $k^\Gamma_\omega(\tau)>0$ for every $\tau>0$ and every real $\omega>0$;
\item $\ct_n>0$ for every $n\ge1$ and every real $\omega>0$, and likewise
 $J_d(n)>0$ for every real $d>0$;
\item $\Kt_\omega$ and $\Kh_\omega$ are completely monotone on $(b,\infty)$ for
 every real $\omega>0$ (Proposition~\ref{prop:cmuncapped});
\item and the only object in the decomposition whose positivity depends on
 $\omega$ at all is $\frac{p+a}{p-a}$, of causal kernel
 $\delta_0+2a\,e^{at}\mathbf 1_{t>0}$, positive exactly when $a\ge0$, that is
 when $m\le2$.
\end{itemize}
Hence the transfer has exactly one positivity threshold in $m$, and it is at
$m=2$.
\end{proposition}

\begin{proof}
$\Gamma(\omega)>0$ and $\sinh\tau>0$ for $\omega,\tau>0$, which gives the first
item from \eqref{eq:kgamma}; each local factor $1-q^{-2\omega}$ and $1-q^{-d}$
lies in $(0,1)$, which gives the second from \eqref{eq:comb} and
Corollary~\ref{cor:oneobject}; the third is
Proposition~\ref{prop:cmuncapped}; the fourth is the causal kernel of
$1+\frac{2a}{p-a}$.
\end{proof}

\begin{remark}[The one other feature at an integer, and it is the wrong
kind]\label{rem:regularity}
At $m=3$, that is $\omega=1$, the archimedean exponent $\omega-1=\frac{m-3}2$
vanishes and $k^\Gamma_\omega$ passes from unbounded at the origin ($m<3$)
through bounded ($m=3$, where it is exactly $2\pi e^{\tau/2}$) to vanishing there
($m>3$). That is a \emph{regularity} threshold, not a positivity one ---
$k^\Gamma_\omega$ is positive throughout --- and it points the wrong way: at
$m=3$ the point count \emph{regains} complete monotonicity, while nothing on the
transfer side is regained or lost.
\end{remark}

\subsection{The two positivity sets diverge}\label{sec:divergence}

\begin{table}[ht]\centering\small
\caption{The two positivity statements on $m>1$. They agree on $(1,2)$ and
differ on $(2,3)$; and at $m=3$ the point count is completely monotone while the
transfer's rational factor is not.}\label{tab:divergence}
\begin{tabular}{@{}llll@{}}
\toprule
$m$ & $\theta(it)^{m}$ compl.\ mon. & $\frac{p+a}{p-a}$ compl.\ mon. & criterion has content\\
\midrule
$(1,2)$ & \textbf{no} --- $r_m(3)<0$ & yes ($a>0$) & \textbf{yes}\\
$2$ & yes (rank-two lattice) & yes ($a=0$, factor $\equiv1$) & no (vacuous)\\
$(2,3)$ & \textbf{no} --- $r_m(7)<0$ & no ($a<0$) & no (vacuous)\\
$3$ & yes (rank-three lattice) & no ($a<0$) & no (vacuous)\\
\bottomrule
\end{tabular}
\end{table}

\begin{corollary}[No equivalence]\label{cor:noequiv}
The point count's failure set contains $(2,3)$, where the criterion is already
vacuous by Proposition~\ref{prop:vacuity}; and at $m=3$ the point count is
completely monotone while $\frac{p+a}{p-a}$ is not. Hence no statement of the
form ``$\theta(it)^{m}$ is completely monotone if and only if [any positivity of
the transfer]'' can hold.
\end{corollary}

\begin{proof}
Propositions~\ref{prop:forced} and~\ref{prop:vacuity} give the first clause and
Propositions~\ref{prop:onethreshold} and~\ref{prop:interp} the second; at $m=3$
the coefficients $r_3(n)$ are representation numbers of $\Z^{3}$ and are
nonnegative, while $a=-\frac12<0$.
\end{proof}

\begin{remark}[The agreement is about the number two]\label{rem:forcedagreement}
The criterion's window is bounded below by $m=1$, where $\omega=0$ and the
transfer degenerates to the identity, and above by $m=2$, where $a$ changes sign.
Those are the two smallest ranks, and the first gap of the rank sequence is
$(1,2)$ by definition. Any pair of thresholds sitting at the two smallest ranks
produces the first gap. The coincidence of intervals is therefore a statement
about the number $2$ and not about either positivity, and it is not evidence of
a common mechanism.
\end{remark}

\subsection{Where the rank sits}\label{sec:whererank}

What does separate the two sides is the \emph{position} the rank occupies in each.

\begin{proposition}[Exponent versus binomial index]\label{prop:position}
Every positivity in the transfer is produced by one of two mechanisms, in each of
which the rank enters as an exponent of a positive quantity.
\begin{itemize}
\item \textbf{Beta densities.} For $\beta>\alpha$, \eqref{eq:betadensity}
 represents $\Gamma(z+\alpha)/\Gamma(z+\beta)$ by the density
 $\Gamma(\beta-\alpha)^{-1}e^{-\alpha t}(1-e^{-t})^{\beta-\alpha-1}$, positive
 for every real $\beta-\alpha>0$. This is the archimedean factor
 $k^\Gamma_\omega\propto(2\sinh\tau)^{\omega-1}e^{\tau/2}$ and both Gamma ratios
 of the regrouping \eqref{eq:regroup}.
\item \textbf{Euler factors.} The comb's local factor at a prime $q$ is
 \begin{equation}\label{eq:eulerlocal}
  \frac{1-q^{-u}}{1-q^{-(u-d)}}
  =1+\sum_{k\ge1}q^{(k-1)d}\big(q^{d}-1\big)\,q^{-ku},
 \end{equation}
 whose coefficients are positive for every real $d>0$ because $q^{d}>1$;
 multiplying over $q\mid n$ recovers $J_d(n)=n^{d}\prod_{q\mid n}(1-q^{-d})$.
\end{itemize}
In both, positivity is a \emph{local} inequality --- at one delay $\tau$, at one
prime $q$ --- in which $d$ appears as an exponent, and an exponent of a positive
quantity is positive at every real value. The point count has neither structure:
by Proposition~\ref{prop:interp} its unique continuation is a binomial series in
which the rank appears as an \emph{index}, and by
Proposition~\ref{prop:signlaw} $\binom mj$ changes sign at $j=\lfloor m\rfloor+2$.
\end{proposition}

\begin{proof}
The Beta integral is the substitution $x=e^{-t}$ in $B(z+\alpha,\beta-\alpha)$.
For \eqref{eq:eulerlocal}, expand
$(1-q^{-u})\sum_{k\ge0}q^{k(d-u)}$; the coefficient of $q^{-ku}$ is
$q^{kd}-q^{(k-1)d}=q^{(k-1)d}(q^{d}-1)$, and multiplicativity over the primes
dividing $n$ gives the Jordan totient. The last clause is
Propositions~\ref{prop:interp} and~\ref{prop:signlaw}.
\end{proof}

\begin{remark}[Reading]\label{rem:reading}
This is what the coincidence of intervals was hiding. The transfer's positivity
was never inherited from a lattice: there is no proof of complete monotonicity
that goes through one, and by Remark~\ref{rem:rankreading} a ratio in which the
lattice cancels could not carry one. It comes from a Beta density and an Euler
product, both of which take the rank as an exponent and therefore continue to
every real value without noticing that a lattice has stopped existing. The point
count takes the rank as a binomial index and notices immediately. The two
positivity statements sit on the same family because the family has one
parameter, and they part company because that parameter occupies a different
position in each.
\end{remark}

\begin{corollary}[A factorization exhibiting the failure does not
exist]\label{cor:nofactorization}
No factor of the decomposition $K_\omega=R_\omega K^\Gamma_\omega K^\zeta_\omega$
loses positivity at any real $d>0$: by
Proposition~\ref{prop:onethreshold} each is positive throughout, and the only
factor that changes character does so at $m=2$, an integer, for a reason internal
to $\Lambda(u)=\Lambda(1-u)$. In particular there is no factorization of the
transfer exhibiting \emph{which} factor loses positivity at fractional $d$,
because none does.
\end{corollary}
